% !TEX program = lualatex
% !TeX encoding = UTF-8
\PassOptionsToPackage{unicode}{hyperref}
\PassOptionsToPackage{bookmarksnumbered,bookmarksopen}{hyperref}
\documentclass[12pt,a4paper]{article}

% ============================================================
% 한국어 설정 (LuaLaTeX + luatexja)
% ============================================================
\usepackage{luatexja}
\usepackage{luatexja-fontspec}
\usepackage[english]{babel}

% 한국어 고딕체 (sans) – Noto Sans KR
\setsansjfont{Noto Sans KR}[
UprightFont = * Regular,
BoldFont = * Bold,
Script = Hangul,
Language = Korean
]

% 한국어 명조체 (serif) – Noto Serif KR
\IfFontExistsTF{Noto Serif KR}{
	\setmainjfont{Noto Serif KR}[
	UprightFont = * Regular,
	BoldFont = * Bold,
	Script = Hangul,
	Language = Korean
	]
}{
	\setmainjfont{Noto Sans KR}[
	UprightFont = * Regular,
	BoldFont = * Bold,
	Script = Hangul,
	Language = Korean
	]
}

% 고정폭 폰트 – 라틴 문자는 Latin Modern Mono, 한글은 Noto Sans KR
\setmonojfont{Noto Sans KR}[
UprightFont = * Regular,
BoldFont = * Bold,
Script = Hangul,
Language = Korean
]

% 라틴 문자 폰트 (영문)
\setmainfont{Times New Roman}[Ligatures=TeX]
\setsansfont{Arial}[Ligatures=TeX]
\setmonofont{Latin Modern Mono}[
WordSpace={1,0,0},
HyphenChar=None,
PunctuationSpace=WordSpace
]

% ============================================================
% 패키지 (원본과 동일)
% ============================================================
\usepackage{amsmath,amssymb,amsfonts,mathtools}
\usepackage{geometry}
\geometry{left=2cm,right=2cm,top=2cm,bottom=2cm}
\usepackage{booktabs}
\usepackage{hyperref}
\usepackage{url}
\usepackage{graphicx}
\usepackage{caption}
\usepackage{subcaption}
\usepackage{float}
\usepackage{siunitx}
\usepackage{bm}
\usepackage{pgfplots}
\pgfplotsset{compat=1.18}
\usepackage{xcolor}
\usepackage{multicol}
\usepackage{tikz}
\usetikzlibrary{positioning, arrows.meta, calc}

% 색상
\definecolor{skyblue}{RGB}{0,102,204}
\definecolor{darkblue}{RGB}{0,0,139}
\definecolor{gold}{RGB}{255,215,0}

% YUCT 매크로
\newcommand{\eps}{\varepsilon}
\newcommand{\kappac}{\kappa_c}
\newcommand{\Keff}{K_{\mathrm{eff}}}
\newcommand{\betaf}{\beta}
\newcommand{\df}{d_{\!f}}
\newcommand{\Mpl}{M_{\mathrm{Pl}}}
\newcommand{\Lpl}{l_{\mathrm{Pl}}}
\newcommand{\tP}{t_{\mathrm{P}}}
\newcommand{\Sodd}{S_{\mathrm{odd}}}
\newcommand{\Seven}{S_{\mathrm{even}}}
\newcommand{\Dtau}{\Delta\tau_{\min}}
\newcommand{\Lzero}{L_0}
\newcommand{\Omegarot}{\Omega_{\mathrm{rot}}}
\newcommand{\Mbar}{M_{\mathrm{bar}}}
\newcommand{\muniv}{M_{\mathrm{univ}}}

\hypersetup{
	colorlinks=true,
	linkcolor=blue,
	citecolor=blue,
	urlcolor=blue,
	pdftitle={YUCT 임계 질량과 우주 계층: 인간 및 AI 독자를 위한 자기 완결적 요약},
	pdfauthor={Alexey V. Yakushev}
}

\title{\textbf{YUCT 임계 질량과 우주 계층: 인간 및 AI 독자를 위한 자기 완결적 요약}}
\author{Alexey V. Yakushev}
\date{2026년 4월}

\begin{document}
	\maketitle
	
	\begin{abstract}
		본 문서는 야쿠셰프 통합 조정 이론(Yakushev Unified Coordination Theory, YUCT)의 국소 빅뱅 임계 질량, 관련 인플레이션 매개변수, 우주의 전체적 회전, 그리고 우주 질량 계층에 관한 주요 결과를 간결하고 자기 완결적으로 요약한다. 모든 방정식은 YUCT 관계식으로 명시되어 있다. 명확한 공식, 표, 반증 가능한 예측을 통해 인간 독자와 AI 언어 모델 모두가 쉽게 이해할 수 있도록 구성되었다. 플랑크 규모에서 관측 가능한 우주까지의 완전한 질량 계층이 제시되어, 단일한 조정 스케일링 법칙을 통해 양자, 항성, 우주론적 규모가 통합됨을 보여준다.
	\end{abstract}
	
	\section{YUCT의 기본 공리}
	
	\begin{itemize}
		\item \textbf{조정 효율} (YPSDC 프로토콜, 부록 B):
		\[
		\Keff = \frac{H(\mathcal{A})}{H(\kappa)}, \qquad\text{(YUCT-1)}
		\]
		여기서 \(H(\mathcal{A})\)는 활성화된 작용의 엔트로피, \(H(\kappa)\)는 전송된 짧은 색인의 엔트로피이다.
		
		\item \textbf{보편적 오차 스케일링 법칙} (부록 L):
		\[
		\varepsilon = \kappac \alpha \Keff^{-\betaf}, \qquad \betaf = \frac{2}{3},\ \kappac = \frac{1}{3}. \qquad\text{(YUCT-2)}
		\]
		
		\item \textbf{대수적 루프} (부록 Y):
		\[
		\Sodd = \frac{6}{5},\quad \Seven = \frac{4}{5},\quad \frac{\Sodd}{\Seven} = \frac{3}{2} = \frac{1}{\betaf}. \qquad\text{(YUCT-3)}
		\]
		기본 길이 및 시간 양자:
		\[
		\Lzero = \frac{2}{3}\Lpl,\qquad \Dtau = \frac{2}{3}\tP. \qquad\text{(YUCT-4)}
		\]
		
		\item \textbf{정보 포화로부터의 프랙탈 차원} (부록 L):
		\[
		d_f = \frac{11}{5}=2.2. \qquad\text{(YUCT-5)}
		\]
	\end{itemize}
	
	\section{국소 빅뱅의 임계 질량}
	
	조정 네트워크의 일관성에서 유도된 스케일링 법칙
	\[
	\Keff(R) = \left(\frac{R}{\Lzero}\right)^{\betaf d_f/2}, \qquad\text{(YUCT-6)}
	\]
	과 임계 조건 \(\varepsilon \approx 1\) (즉, \(\Keff = K_{\mathrm{crit}} = 3^{3/2}\))을 사용하면, 중력적으로 속박된 전(前)붕괴 영역의 임계 질량이 얻어진다:
	\[
	M_{\mathrm{crit}} = 3 M_{\mathrm{Pl}} \approx 6.5\times 10^{-8}\ \mathrm{kg}. \qquad\text{(YUCT-7)}
	\]
	이 유도에서는 슈바르츠실트 반지름 \(R_S = 2GM/c^2\)과 대수적 루프 \(\Lzero = \frac{2}{3}\Lpl\)이 사용된다. 양자 보정은 선도 차수에서 계수 3을 변화시키지 않는다.
	
	\section{조정 상전이로서의 우주론적 인플레이션}
	
	인플레이션 관측량은 동일한 공리로부터 예측된다:
	\begin{align}
		n_s &= 1 - 2\betaf^2 + \frac{4}{3}\betaf^3 + \cdots \approx 0.965, \qquad\text{(YUCT-8)}\\
		r &= 8(2\betaf)^2 = 32\betaf^2 \approx 0.07, \qquad\text{(YUCT-9)}\\
		f_{\mathrm{NL}}^{\mathrm{local}} &= \frac{5}{3}\betaf \approx 1.12. \qquad\text{(YUCT-10)}
	\end{align}
	이러한 예측은 CMB‑S4, LiteBIRD, Euclid, SPHEREx에 의해 반증 가능하다.
	
	\section{우주의 전체적 회전}
	
	CMB 쌍극자의 일부는 각속도 \(\omega\)를 갖는 전체적인 강체 회전에 기인한다. 쌍극자 진폭은 적색편이 \(z\)와 함께 다음과 같이 증가한다:
	\[
	v_{\mathrm{dip}}(z) = v_{\mathrm{local}} + \omega\, r(z). \qquad\text{(YUCT-11)}
	\]
	NVSS, Quaia 및 다른 서베이에서 관측된 증가로부터,
	\[
	\omega \approx 8.5\times 10^{-22}\ \mathrm{rad/s},\qquad
	T_{\mathrm{rot}} = \frac{2\pi}{\omega} \approx 2.3\times 10^{14}\ \mathrm{yr}. \qquad\text{(YUCT-12)}
	\]
	무차원 회전 매개변수는
	\[
	\Omega_{\mathrm{rot}} \equiv \frac{\omega}{H_0} = \frac{\Seven}{\Sodd} \betaf^2 \mathcal{F}(d_f)\left(\frac{L_0}{R_H}\right)^{2\betaf} \approx 3.9\times 10^{-4}, \qquad\text{(YUCT-13)}
	\]
	이며, 허블 매개변수는 \(H_0 = \betaf\,\omega\)를 만족한다. 우주의 바리온 질량은 회전으로부터 다음과 같이 유도된다:
	\[
	\Mbar = \Omega_b\frac{\betaf^2 c^3}{G H_0} \approx (2.3\pm 0.2)\times 10^{51}\ \mathrm{kg}. \qquad\text{(YUCT-14)}
	\]
	
	\section{양자에서 우주까지의 완전한 질량 계층}
	
	YUCT에서 플랑크 규모에서 관측 가능한 우주까지의 모든 질량은 단일한 조정 스케일링 법칙을 따른다. 표~\ref{tab:full-hierarchy}는 계층 수준, 전형적인 질량, 예시 및 해당하는 YUCT 스케일링 관계를 나열한다. 그림~\ref{fig:hierarchy-scheme}은 이산적인 사다리 단계를 시각화한다.
	
	\begin{table}[htbp]
		\centering
		\caption{양자에서 우주까지의 완전한 질량 계층 (YUCT).}
		\label{tab:full-hierarchy}
		\small
		\begin{tabular}{l c c l}
			\toprule
			\textbf{수준} & \textbf{전형적 질량 (kg)} & \textbf{예시} & \textbf{YUCT 스케일링 관계} \\
			\midrule
			플랑크 씨앗 & \(M_{\mathrm{crit}} = 3M_{\mathrm{Pl}} \approx 6.5\times10^{-8}\) & 국소 빅뱅 씨앗 & \(M_{\mathrm{crit}} = 3 M_{\mathrm{Pl}}\) (YUCT-7) \\
			전자 & \(m_e \approx 9.1\times10^{-31}\) & 가장 가벼운 하전 페르미온 & \(m_e = M_{\mathrm{Pl}} (2/3)^{96+\sigma}\xi_e\) \\
			양성자 & \(m_p \approx 1.67\times10^{-27}\) & 안정한 바리온 & \(m_p = M_{\mathrm{Pl}} (2/3)^{96+\sigma}\xi_p\) \\
			백색왜성 & \(\sim 10^{30}\) (\(0.6M_\odot\)) & 찬드라세카르 한계 & \(M_{\mathrm{WD}} \propto M_{\mathrm{Pl}}^{3}/m_p^2\) \\
			중성자별 & \(\sim 3\times10^{30}\) (\(1.4M_\odot\)) & TOV 한계 & \(M_{\mathrm{NS}} \propto M_{\mathrm{Pl}}^{3}/m_p^2\) \\
			초대질량 블랙홀 & \(\sim 10^{36}\)–\(10^{40}\) & 은하 중심 & \(M_{\mathrm{BH}} \propto M_{\mathrm{Pl}} (2/3)^{N}\) \\
			관측 가능한 우주 & \(M_{\mathrm{univ}} \approx 1.5\times10^{53}\) & \(R_H\) 내의 총 질량 & \(M_{\mathrm{univ}} = \frac{\betaf^2 c^3}{G H_0}\) (YUCT-14) \\
			바리온 우주 & \(M_{\mathrm{bar}} \approx 2.3\times10^{51}\) & 가시 물질 & \(M_{\mathrm{bar}}/m_p = (M_{\mathrm{Pl}}/m_e)^{3.5}\) (YUCT-15) \\
			\bottomrule
		\end{tabular}
	\end{table}
	
	\begin{figure}[htbp]
		\centering
		\begin{tikzpicture}[
			level/.style={rectangle, draw=darkblue!70, fill=skyblue!10, thick, minimum width=7cm, minimum height=0.7cm, rounded corners, align=center, font=\small},
			arrow/.style={->, >=Stealth, thick, color=darkblue},
			]
			\node[level] (planck) {플랑크 씨앗: \(M_{\mathrm{crit}} = 3M_{\mathrm{Pl}}\)};
			\node[level, below=0.6cm of planck] (electron) {전자: \(m_e\)};
			\node[level, below=0.6cm of electron] (proton) {양성자: \(m_p\)};
			\node[level, below=0.6cm of proton] (wd) {백색왜성: \(\sim 0.6 M_\odot\)};
			\node[level, below=0.6cm of wd] (ns) {중성자별: \(\sim 1.4 M_\odot\)};
			\node[level, below=0.6cm of ns] (bh) {초대질량 블랙홀: \(10^6\)–\(10^{10} M_\odot\)};
			\node[level, below=0.6cm of bh] (univ) {관측 가능한 우주: \(M_{\mathrm{univ}}\)};
			\draw[arrow] (planck) -- (electron);
			\draw[arrow] (electron) -- (proton);
			\draw[arrow] (proton) -- (wd);
			\draw[arrow] (wd) -- (ns);
			\draw[arrow] (ns) -- (bh);
			\draw[arrow] (bh) -- (univ);
		\end{tikzpicture}
		\caption{YUCT에서 플랑크 씨앗에서 우주까지의 이산적 질량 사다리. 각 단계는 조정 장의 공명에 해당한다.}
		\label{fig:hierarchy-scheme}
	\end{figure}
	
	우주 질량 계층은 더욱이 다음 대수적 관계로 요약된다:
	\[
	\boxed{\frac{\Mbar}{m_p} = \left(\frac{M_{\mathrm{Pl}}}{m_e}\right)^{\mathcal{S}}},\qquad
	\mathcal{S} \equiv \Sodd + \Seven + \frac{\Sodd}{\Seven} = 3.5. \qquad\text{(YUCT-15)}
	\]
	수치적으로 \((M_{\mathrm{Pl}}/m_e)^{3.5} \approx 1.88\times 10^{78}\)이며, 이는 관측값 \(\Mbar/m_p \approx 1.4\times 10^{78}\)과 1.3배 이내로 일치한다. 그 차이는 재가열 효율 \(\epsilon_{\mathrm{reh}}\approx 0.85\)와 기하학적 인자 \(\mathcal{F}_{\mathrm{rot}}\approx 0.9\)로 설명된다. 이는 보편적 오차 지수를 지배하는 동일한 대수적 상수 \(S_{\mathrm{odd}}, S_{\mathrm{even}}\)이 우주의 바리온 예산까지도 결정함을 보여준다.
	
	\section{우주 천체의 계층과 임계 회전}
	
	축퇴압으로 지지되는 자기 중력계에 대해, YUCT는 임계 각속도의 보편적 스케일링을 예측한다:
	\[
	\Omega_{\mathrm{crit}} \propto M^{\betaf},\qquad \betaf = 2/3. \qquad\text{(YUCT-16)}
	\]
	표~\ref{tab:objects}와 그림~\ref{fig:omega-mass}는 행성에서 관측 가능한 우주까지의 이 계층을 예시한다.
	
	\begin{table}[htbp]
		\centering
		\caption{선택된 우주 천체의 임계 매개변수 (YUCT).}
		\label{tab:objects}
		\small
		\begin{tabular}{l c c c}
			\toprule
			\textbf{천체} & \textbf{질량 (\(M_\odot\))} & \(\log_{10}(\Omega_{\mathrm{crit}}/\mathrm{s^{-1}})\) & \textbf{스케일링 법칙} \\
			\midrule
			가스 거대 행성 (목성) & 0.001 & \(-3.8\) & \(R\sim M^{-0.67}\) \\
			적색 왜성 (M5V) & 0.2 & \(-4.5\) & \(L\sim M^{2.3}\) \\
			백색왜성 & 0.6 & 0.2 & \(R\sim M^{-1/3}\) \\
			중성자별 & 1.4 & 3.7 & \(M\sim\Lambda^{0.67}\) \\
			초질량 중성자별 & 2.2 & 4.2 & \(\Omega_{\mathrm{crit}}\sim M^{0.67}\) \\
			항성 질량 블랙홀 & 5 & -- & \(R_S\sim M\) \\
			관측 가능한 우주 & \(7.5\times10^{22}\) & \(-17.7\) (\(H_0\)) & \(\rho_c\sim H^2\sim K_{\mathrm{eff}}^{2/3}\) \\
			\bottomrule
		\end{tabular}
	\end{table}
	
	\begin{figure}[htbp]
		\centering
		\begin{tikzpicture}
			\begin{loglogaxis}[
				xlabel={질량 \(M\) (\(M_\odot\))}, ylabel={\(\Omega_{\mathrm{crit}}\) (s\(^{-1}\))},
				grid=both, width=0.9\textwidth, height=0.6\textwidth,
				legend pos=north west, title={밀집 천체의 임계 회전 (YUCT)},
				xmin=0.08, xmax=15, ymin=5e-2, ymax=2e4,
				]
				\addplot[domain=0.1:10, samples=2, thick, blue, dashed] {10^(0.67*log10(x)+0.5)};
				\addlegendentry{\(\Omega_{\mathrm{crit}}\propto M^{2/3}\) (백색왜성)};
				\addplot[domain=1:3, samples=2, thick, green!60!black, dashed] {10^(0.67*log10(x)+3.7)};
				\addlegendentry{\(\Omega_{\mathrm{crit}}\propto M^{2/3}\) (중성자별)};
				\addplot[only marks, mark=*, mark size=3pt, red] coordinates {(0.6,1.6)};
				\addplot[only marks, mark=square*, mark size=3pt, green!60!black] coordinates {(1.4,5000) (2.2,15000)};
				\addplot[only marks, mark=triangle*, mark size=3pt, orange] coordinates {(0.2,3e-5)};
			\end{loglogaxis}
		\end{tikzpicture}
		\caption{임계 각속도 대 질량 (YUCT). 기울기는 \(\beta=2/3\)을 확인한다.}
		\label{fig:omega-mass}
	\end{figure}
	
	\section{반증 가능한 예측}
	
	\begin{enumerate}
		\item \textbf{원시 블랙홀 유물:} 암흑 물질의 \(f_{\mathrm{relic}} \approx 0.05\). 만약 \(f_{\mathrm{relic}}>0.2\) 또는 \(<0.01\) (95\% CL)이면 반증됨.
		\item \textbf{인플레이션 매개변수:} \(n_s = 0.965\pm0.005\), \(r\approx 0.07\). \(r<0.01\) 또는 \(n_s\)가 \(0.960\pm0.010\) 범위 밖이면 배제됨.
		\item \textbf{비가우시안성:} \(f_{\mathrm{NL}}^{\mathrm{local}} \approx 1.12\). 만약 측정값이 \(>3\sigma\)에서 \(<0.5\) 또는 \(>1.8\)이면 반증됨.
		\item \textbf{\(E=mc^2\)로부터의 실험실 편차:} \(\Delta E/E \sim \kappa_c\alpha K_{\mathrm{eff}}^{-\beta}\). \(10^{-18}\) 감도에서 귀무 결과는 YUCT에 대한 도전이 됨.
		\item \textbf{우주 회전 주기:} \(T_{\mathrm{rot}} \approx 2.3\times10^{14}\) 년. 미래의 서베이(Euclid, SKA)가 \(\omega\)를 정밀화할 것임.
	\end{enumerate}
	
	\section{결론}
	
	YUCT의 공리는 국소 빅뱅의 임계 질량, 인플레이션 관측량, 우주의 전체적 회전, 그리고 플랑크 씨앗에서 관측 가능한 우주까지의 완전한 질량 계층을 통합한다. 모든 예측은 정량적이며 반증 가능하다. 위에 핵심 방정식들이 요약되어 있으며, 각각 YUCT 관계식으로 표시되었다. 상세한 유도는 인용된 부록들에서 찾을 수 있다.
	
	\section*{데이터 이용 가능성}
	
	전체 유도 및 코드: \url{https://github.com/Alexey-Yakushev-YUCT/YPSDC}
	
	\section*{YUCT 부록 참고문헌}
	\begin{itemize}
		\item 부록 L: 프랙탈 조정 오차 스케일링 – 보편적 법칙 \(\beta=2/3\), \(d_f=11/5\).
		\item 부록 Y: 수학적 기초와 KPZ 유도.
		\item 부록 P: 중력의 온도 의존성.
		\item 부록 M: 조정 상전이로서의 인플레이션.
		\item 부록 \(\Omega\): 우주의 전체적 회전.
		\item 부록 B: YPSDC 프로토콜.
	\end{itemize}
	
\end{document}